% ECONOMICS_PROBLEM: {"id": "C73-71", "title": "Non-equilibrium learning attractors", "jel_code": "C73", "source_id": "OP-0284"}
% !TeX program = xelatex
\documentclass[11pt,a4paper]{article}
\usepackage[margin=23mm]{geometry}
\usepackage{fontspec}
\setmainfont{Latin Modern Roman}
\usepackage{amsmath,amssymb,mathtools}
\usepackage{xcolor,enumitem,booktabs,longtable,array,xurl,hyperref}
\definecolor{muted}{HTML}{53636C}
\hypersetup{colorlinks=true,urlcolor=blue,linkcolor=blue}
\setlength{\parindent}{0pt}
\setlength{\parskip}{5pt}
\newcommand{\R}{\mathbb R}
\newcommand{\E}{\mathbb E}
\newcommand{\N}{\mathbb N}
\newcommand{\ind}{\mathbf 1}
\newcommand{\Var}{\operatorname{Var}}
\newcommand{\Cov}{\operatorname{Cov}}
\newcommand{\supp}{\operatorname{supp}}
\DeclareMathOperator*{\argmin}{arg\,min}
\DeclareMathOperator*{\argmax}{arg\,max}
\newcommand{\entrymeta}[3]{{\sffamily\small\color{muted}\textbf{\texttt{#1}}\quad|\quad #2\par #3\par}}
\newcommand{\Problemref}[1]{\texttt{#1}}
\newcommand{\JELref}[2]{\texttt{#2} (JEL \texttt{#1})}
\begin{document}
\section*{Non-equilibrium learning attractors}
\textbf{Problem ID:} \texttt{C73-71}\quad \textbf{JEL:} \texttt{C73}\par
% BEGIN ECONOMICS STATEMENT
\label{problem:OP-0284}
{\sffamily\small\color{muted}Original subject group: Game theory and strategic interaction\par}
\label{definitions:problem:30007566}
\textbf{Audit classification:} Explicit published open question. The v4 section 4 Algorithmic problems item 1 explicitly asks for a local-source converse. The revised statement fixes an attraction-based interpretation, with the uniform convention supported by Biggar--Shames (2023); Lyapunov stability alone is distinguished.
\subsection*{Definitions and precise research target}
\textbf{Model and research target.} Let a finite two-player game have rational payoff matrices \(A,B\in\mathbb Q^{r\times s}\), mixed-strategy space \(X=\Delta_r\times\Delta_s\), and replicator flow \(\varphi_t\) defined by
\[
 \dot x_a=x_a((Ay)_a-x^TAy),\qquad
 \dot y_b=y_b((x^TB)_b-x^TBy).
\]
Here \(\Delta_k\) is the probability simplex on \(k\) actions. The preference graph has pure profiles as vertices and an arc for each unilateral weak payoff improvement. Let \(H\) be a strongly connected component with no outgoing arcs. Its content \(K_H\) consists of mixed profiles whose product distribution is supported entirely on \(H\). A compact invariant set \(S\subseteq X\) is attracting if some positively invariant neighborhood \(U\) of \(S\) satisfies \(\sup_{z\in U}\operatorname{dist}(\varphi_t(z),S)\to0\); it is an attractor if it has no proper compact invariant attracting subset. Apply these definitions to \(S=K_H\). This is the attraction-based interpretation of the source question; Lyapunov stability alone means that every sufficiently close initial condition stays close for all future times and does not imply attraction. A subgame \(Y\) is a product of faces of the two simplexes. A local source of \(H\) is a point \(z\in K_H\cap Y\), with \(Y\not\subseteq K_H\), that is a quasi-strict Nash equilibrium of the negated-payoff game restricted to \(Y\): its best responses are exactly its supported actions. Determine whether every sink component whose content fails to be an attractor necessarily has such a local source in some subgame. Prove the implication for all finite bimatrix games in the stated domain, or supply a counterexample and identify the alternative obstruction.

\emph{Definitions and maintained assumptions (editorial unless a source definition is named).} Use \(\Delta_r=\{x\ge0:\sum_a x_a=1\}\). The preference graph has an arc \((a,b)\to(a',b)\) when \(A_{a'b}\ge A_{ab}\), and similarly for player two; endpoints are distinct. A sink is a strongly connected component with no outgoing arc. Define \(K_H=\{(x,y):\operatorname{supp}(x)\times\operatorname{supp}(y)\subseteq H\}\). The polynomial replicator field defines a global flow on compact \(X\). To make the audit’s convention explicit, an attracting set is nonempty compact invariant \(S\) with a relative positively invariant neighborhood \(U\) and \(\sup_{z\in U}\operatorname{dist}(\varphi_t(z),S)\to0\). Allow \(S=U=X\), otherwise requiring a strictly larger neighborhood would contradict the source’s whole-space examples. An attractor is minimal among such attracting sets. A subgame is a product of faces. Quasi-strict Nash in the negated subgame means every supported action minimizes its original expected payoff, all supported actions tie, and each unsupported action in that subgame has strictly larger original payoff. The finite rational bimatrix restriction is a specialization of the source’s broader domain. For a sink content, attracting implies minimal attracting: an attracting subset contains a pure vertex of \(H\), hence all of \(H\); the chain-component structure then forces it to contain \(K_H\). This uses the 2023 reference's Lemma~3.5 and Theorem~5.2, together with the fact that an attracting set that meets a chain component contains that component. Thus minimality does not introduce a distinct sink-content obstruction.
\subsection*{Variants and assumption changes}
\begin{enumerate}
\item \textbf{Many players (source question, editorial formalization).} For \(N\) finite players replace product supports and the replicator equations accordingly. Ask whether every nonattracting sink content admits a local source in some product-face subgame.
\item \textbf{Weaker convergence (source research direction).} Require only \(\operatorname{dist}(\varphi_t(z),K_H)\to0\) for every interior \(z\), without uniformity or attraction of boundary histories. Determine the appropriate witness criterion; this is distinct from the preceding minimal-attractor question.
\end{enumerate}
\subsection*{References and source audit}
\begin{enumerate}
\item Definition2.1; §4 Algorithmic problems item1, printed p.12. Explicit attraction-converse question checked; the formal statement uses attraction rather than interpreting every occurrence of ``stable'' as Lyapunov stability. Added domain restrictions are editorial. The uniform-attraction convention has explicit support in the 2023 reference. \url{https://arxiv.org/pdf/2502.07975v3}
\item Definition2.1; §4 Algorithmic problems item1, printed p.12. Explicit attraction-converse question checked; the formal statement uses attraction rather than interpreting every occurrence of ``stable'' as Lyapunov stability. Added domain restrictions are editorial. The uniform-attraction convention has explicit support in the 2023 reference. \url{https://arxiv.org/pdf/2502.07975v4}
\item Biggar--Shames (2023), Definition3.2 verifies the uniform-attraction convention; Lemma3.5 and Theorem5.2 give the sink/chain structure. \url{https://proceedings.mlr.press/v201/biggar23a/biggar23a.pdf}
\end{enumerate}
\begin{enumerate}\raggedright
\item\hypertarget{definitions:source:30007566:1}{} Oliver Biggar; Christos H. Papadimitriou. 2025. \emph{Sink equilibria and the attractors of learning in games}. Definition2.1; §4 Algorithmic problems, item1, printed p.12, version3 (26October2025).
\url{https://arxiv.org/pdf/2502.07975v3}.
\item\hypertarget{definitions:source:30007566:2}{} Oliver Biggar; Christos H. Papadimitriou. 2026. \emph{Sink equilibria and the attractors of learning in games}. §4 Algorithmic problems, item1, printed p.12, version4 (4March2026).
\url{https://arxiv.org/pdf/2502.07975v4}.
\item\hypertarget{definitions:source:30007566:3}{} Oliver Biggar; Iman Shames. 2023. \emph{The Replicator Dynamic, Chain Components and the Response Graph}. Definition3.2, Lemma3.5, Theorem5.2; Proceedings of Machine Learning Research201,1–22.
\url{https://proceedings.mlr.press/v201/biggar23a/biggar23a.pdf}.
\end{enumerate}

\subsection*{Common conventions: Source mathematical conventions}

These conventions apply to each entry unless explicitly replaced.

\textbf{Models and identification.} A primitive model \(m\) specifies all probability laws, feasible choices, information, equilibrium and measurement rules used by its target. The admissible class \(\mathcal M\) is fixed independently of outcomes. If \(P_O(m)\) is its observable probability law and \(\tau(m)\) the target, its identified set at observable law \(P\) is
\[
 \mathcal I_\tau(P)=\{\tau(m):m\in\mathcal M,\ P_O(m)=P\}.
\]
Point identification means that this set is a singleton. An empty set signals incompatibility of data and assumptions. Coordinate bounds need not describe the joint set. Bounds are sharp only if every retained target is attainable or arbitrarily approachable by compatible models. Finite bounds require explicit restrictions on unobserved quantities; unrestricted confounding, hidden wealth, extrapolation or arbitrary equilibrium selection can yield uninformative sets.

\textbf{Causal interventions.} A potential outcome \(Y_i(p)\) is the outcome under a complete policy regime \(p\), including timing, financing, information and market spillovers. The target population and units are fixed before treatment. With interference, use the complete assignment vector or a justified exposure mapping; individual no-interference notation alone cannot identify an economy-wide intervention. A difference of mean potential outcomes does not require the cross-world joint distribution. Conditioning on post-treatment migration, survival, enrollment or employment changes the estimand and needs separate justification.

\textbf{Statistical statements.} An estimator, confidence set, forecast or test specifies a sampling law, model class and loss/coverage criterion. Uniform \((1-\alpha)\)-coverage means \(\inf_{m\in\mathcal M}\Pr_m\{\tau(m)\in C_n\}\geq1-\alpha\), or an explicitly asymptotic version. The whole parameter space trivially has coverage, so informativeness requires a stated width, risk or power objective. A forecasting improvement is relative to a named benchmark, information set and evaluation population.

\textbf{Equilibrium and optimization.} An equilibrium comprises feasible plans, optimality, beliefs where needed, budget constraints and all market/resource clearing. The primitives or a stated benchmark define each correspondence \(\mathcal E(p;m)\). Nonemptiness is assumed or proved, never inferred just from compact policy sets. A selected-equilibrium target uses a declared selection rule; a robust target quantifies over all permitted equilibria. Use suprema or infima unless attainment is established. An \(\varepsilon\)-optimal policy is feasible and within \(\varepsilon\) of the finite optimum value. Report an empty feasible set separately.

\textbf{Welfare and accounting.} Normative utility, interpersonal normalization, social weights, numeraire, discounting and the use of public receipts are primitives. Behavioral utility need not coincide with the welfare criterion. Household welfare includes ownership income and rents once; adding profits again to compensating variation generally double-counts them. A monetary equivalent is defined by an explicit transfer and price convention, with monotonicity and range conditions. Average effects, mechanism contributions, transfers and welfare losses are different targets. Infinite sums require convergence and dynamic feasible plans require terminal or no-Ponzi conditions.

\textbf{Source versions.} A working-paper date, revision date and journal publication year are distinguished. A DOI for a journal version is not automatically the DOI of an earlier manuscript. Locators refer to the stated version and printed pages where specified. A metadata-only check does not verify a claim about a full-text conclusion. Every entry's source subsection narrows the provenance claim accordingly.
% END ECONOMICS STATEMENT
\end{document}
